% TOP_PROBLEM: {"id": 2100403, "title": "Open Problems in Integrable Systems — Around the Birkhoff conjecture", "category_id": 11, "category": {"id": 11, "name": "dynamical_systems", "display_name": "Dynamical Systems", "description": "Problems about long-term behavior of deterministic systems, Hamiltonian dynamics, and periodic orbits.", "slug": "dynamical-systems", "order_index": 11, "created_at": "2026-07-31T15:26:25.671Z"}, "status": "open"}
% FIRST_POSTED: null
% ENTRY_KIND: statement_only
% Imported from: batch08.tex; UnsolvedMath ID 2100403
% Compile twice: lualatex 2100403.tex
\documentclass[11pt]{article}
\usepackage{fontspec}
\setmainfont{Latin Modern Roman}
\usepackage{amsmath,amssymb,mathtools,unicode-math}
\setmathfont{Latin Modern Math}
\usepackage{xurl,hyperref}
\hypersetup{hidelinks,pdftitle={UnsolvedMath Problem 2100403}}
\newcommand{\sref}[2]{\hyperlink{source:#1:#2}{[S#2]}}
\newcommand{\cataloguescope}[1]{\paragraph{Mathematical question.}#1}
\begin{document}
% UnsolvedMath ID 2100403; source catalogue code AMR-020-0403
\setcounter{section}{2100402}
\section{Open Problems in Integrable Systems — Around the Birkhoff conjecture}\label{2100403}
{\small\sffamily UnsolvedMath ID 2100403 · Dynamical Systems}
\subsection{Definitions and mathematical statement}

\paragraph{Shared notation.}$\mathbb N=\{1,2,\ldots\}$, and $\mathbb Z,\mathbb Q,\mathbb R,\mathbb C$ denote the integers, rationals, reals and complex numbers. Any other conventions are specified locally.


Let $\gamma\subset\mathbb{RP}^2$ be a smooth closed strictly convex curve in an affine chart, with a smooth exterior foliation $\mathcal F$ by closed convex curves including $\gamma$ as a leaf. For each tangent line $L$ to $\gamma$, pairing the two nearby intersections of $L$ with the same leaf defines a local involution, fixing the tangency point. Assume this involution is the restriction of a projective transformation of $L$. A first integral is a nonconstant function whose level curves are the foliation leaves. A rational integral is rational in affine coordinates; an algebraic integral is a local branch of a function satisfying a nonzero polynomial equation over the rational-function field. A pencil of conics consists of equations $aQ_0+bQ_1=0$ for fixed homogeneous quadratic forms $Q_0,Q_1$.
\paragraph{Statement recorded in the supplied source.}
Prove that $\gamma$ is a conic (an ellipse in the affine chart) and that $\mathcal F$ consists of conics in one pencil when the foliation admits (i) a rational first integral or (ii) an algebraic first integral. These are the two specified special cases of the source's projective-involution conjecture. \sref{2100403}{2}
\subsection{Short English statement}
Prove the conic-pencil conclusion for a convex exterior foliation with projective tangent-line involutions and a rational or algebraic first integral.
\subsection{Sources}
\begin{description}
\item[\hypertarget{source:2100403:1}{S1}] ULAM.AI and the UnsolvedMath Contributors, UnsolvedMath: A Curated Collection of Open Mathematics Problems, record 2100403 (AMR-020-0403). CC BY 4.0. \url{https://huggingface.co/datasets/ulamai/UnsolvedMath}.
\item[\hypertarget{source:2100403:2}{S2}] A. Bolsinov, V. S. Matveev, E. Miranda and S. Tabachnikov, Open Problems, Questions, and Challenges in Finite-Dimensional Integrable Systems (2018), arXiv:1804.03737, Problem 4.5 and Conjecture 4.3. \url{https://arxiv.org/abs/1804.03737}.
\end{description}
\label{end:2100403}
\end{document}
